% Unit 040 terjemahan Bahasa Indonesia dari chapter3.tex baris 1587--1709.
% Sumber: Methods of Algebra, Volume 2, commit
% 9a5803ff77dd3257484cb177f851a73770a59dd3 (CC BY 4.0).
% stable-unit-id: o014.aljabr2.chapter3.truncation-functors
% Cakupan: Bagian 3.9 lengkap.

% segment-id: o014.li2.chapter3.truncation-functors.h001
\section{Funktor pemenggalan}\label{sec:truncation-functors}
\hypertarget{o014.aljabr2.chapter3.truncation-functors}{ }

% segment-id: o014.li2.chapter3.truncation-functors.p001
Pada awal bagian ini, misalkan $\mathcal{A}$ suatu kategori aditif.

% segment-id: o014.li2.chapter3.truncation-functors.de001
\begin{definisi}\label{def:cplx-cat-variant}
	\index{kompleks!terbatas atas, terbatas bawah, terbatas (bounded above, bounded below, bounded)}
	\index[sym1]{CbA@$\cate{C}^+(\mathcal{A})$, $\cate{C}^-(\mathcal{A})$, $\cate{C}^{\bdd}(\mathcal{A})$}
	\index[sym1]{CstA@$\cate{C}^{[s,t]}(\mathcal{A})$, $\cate{C}^{\geq s}(\mathcal{A})$, $\cate{C}^{\leq t}(\mathcal{A})$}
	Untuk kompleks $X \in \Obj(\cate{C}(\mathcal{A}))$, gunakan istilah
	berikut beserta subkategori penuh yang bersesuaian:
% segment-id: o014.li2.chapter3.truncation-functors.q001
	\begin{center}\small\setlength{\tabcolsep}{3pt}\begin{tabular}{|c|c|c|c|}\hline
		Istilah & terbatas & terbatas bawah & terbatas atas \\
		Syarat & $|n| \gg 0 \implies X^n = 0 $ & $n \ll 0 \implies X^n = 0$ & $n \gg 0 \implies X^n = 0$ \\
		Subkategori penuh & $\cate{C}^{\bdd}(\mathcal{A})$ & $\cate{C}^+(\mathcal{A})$ & $\cate{C}^-(\mathcal{A})$ \\ \hline
	\end{tabular}\end{center}
	
	Secara lebih umum, untuk $-\infty \leq s \leq t \leq +\infty$,
	gunakan $\cate{C}^{[s,t]}(\mathcal{A})$ untuk subkategori penuh yang
	terdiri atas kompleks-kompleks yang memenuhi
	$n \notin [s, t] \implies X^n = 0$, dan definisikan
% segment-id: o014.li2.chapter3.truncation-functors.q002
	\[ \cate{C}^{\geq s}(\mathcal{A}) := \cate{C}^{[s, +\infty]}(\mathcal{A}), \quad \cate{C}^{\leq t}(\mathcal{A}) := \cate{C}^{[-\infty, t]}(\mathcal{A}). \]
\end{definisi}

% segment-id: o014.li2.chapter3.truncation-functors.p002
Semua subkategori penuh tersebut merupakan kategori aditif. Jika
$\mathcal{A}$ merupakan kategori abelian, maka $\cate{C}^+(\mathcal{A})$,
$\cate{C}^-(\mathcal{A})$, dan
$\cate{C}^{\bdd}(\mathcal{A}) = \cate{C}^+(\mathcal{A}) \cap
\cate{C}^-(\mathcal{A})$ merupakan subkategori abelian.\footnote{Catatan
penerjemah (O014-C050): sumber menyatakan bagian kedua tanpa mensyaratkan
$\mathcal{A}$ abelian, padahal pada titik ini $\mathcal{A}$ baru diasumsikan
aditif. Klaim tersebut dibuat bersyarat di sini.} Funktor pergeseran
$[n]$ mempertahankan $\cate{C}^+(\mathcal{A})$,
$\cate{C}^-(\mathcal{A})$, dan $\cate{C}^{\bdd}(\mathcal{A})$, tetapi
memetakan $\cate{C}^{[a,b]}(\mathcal{A})$ ke
$\cate{C}^{[a-n, b-n]}(\mathcal{A})$.

% segment-id: o014.li2.chapter3.truncation-functors.p003
Selain itu, $\mathcal{A}$ secara alami diidentifikasikan dengan
$\cate{C}^{\geq 0}(\mathcal{A}) \cap \cate{C}^{\leq 0}(\mathcal{A})$.

% segment-id: o014.li2.chapter3.truncation-functors.p004
Untuk setiap $\star \in \{+, -, \bdd\}$, konsep homotopi antarmorfisme
(Definisi~\ref{def:homotopy}) dapat dibatasi pada
$\cate{C}^{\star}(\mathcal{A})$. Dengan demikian,
Definisi~\ref{def:cplx-homotopy-cat} mempunyai perumuman berikut.

% segment-id: o014.li2.chapter3.truncation-functors.de002
\begin{definisi}\label{def:cplx-homotopy-cat-variant}
	\index[sym1]{KbA@$\cate{K}^+(\mathcal{A})$, $\cate{K}^-(\mathcal{A})$, $\cate{K}^{\bdd}(\mathcal{A})$}
	Untuk $\star \in \{+, -, \bdd\}$, $\cate{K}(\mathcal{A})$ mempunyai
	subkategori penuh aditif $\cate{K}^\star(\mathcal{A})$ yang tertutup
	terhadap funktor pergeseran.
\end{definisi}

% segment-id: o014.li2.chapter3.truncation-functors.p005
Mulai sekarang, kita mensyaratkan bahwa $\mathcal{A}$ merupakan kategori
abelian. Kita akan mengkaji cara memenggal bagian kompleks $X$ yang berderajat $< n$ (atau
$> n$). Gagasan naif ialah mengganti semua suku pada bagian tersebut dengan
$0$. Walaupun sederhana, cara ini mengacaukan kohomologi; karena itu, kita
memperkenalkan versi yang lebih halus.

% segment-id: o014.li2.chapter3.truncation-functors.de003
\begin{definisi}[Funktor pemenggalan]\label{def:truncation-cplx}
	\index{funktor pemenggalan (truncation functor)}
	\index[sym1]{taun@$\tau^{\leq n}$, $\tau^{\geq n}$}
	\index[sym1]{tautilden@$\tilde{\tau}^{\leq n}$, $\tilde{\tau}^{\geq n}$}
	Misalkan $\mathcal{A}$ kategori abelian dan $n \in \Z$. Untuk suatu
	kompleks $X$, definisikan
% segment-id: o014.li2.chapter3.truncation-functors.g001
	\[\begin{tikzcd}
		[column sep=small, every cell/.append style = {font = \small}]
		\tau^{\leq n} X \arrow[phantom, r, "" {name=A, yshift=2ex}] \arrow[hookrightarrow, d] & \cdots \arrow[r] & X^{n-2} \arrow[r] & X^{n-1} \arrow[r] & \Ker(d^n) \arrow[r] \arrow[hookrightarrow, d] & 0 \arrow[r] & 0 \arrow[r] & \cdots \\
		\tilde{\tau}^{\leq n} X \arrow[hookrightarrow, d] & \cdots \arrow[r] & X^{n-2} \arrow[r] & X^{n-1} \arrow[r] & X^n \arrow[r] & \Coim(d^n) \arrow[r] \arrow[hookrightarrow, d] & 0 \arrow[r] & \cdots \\
		X \arrow[twoheadrightarrow, d] & \cdots \arrow[r] & X^{n-2} \arrow[r] & X^{n-1} \arrow[r] \arrow[twoheadrightarrow, d] & X^n \arrow[r] & X^{n+1} \arrow[r] & X^{n+2} \arrow[r] & \cdots \\
		\tilde{\tau}^{\geq n} X \arrow[twoheadrightarrow, d] & \cdots \arrow[r] & 0 \arrow[r] & \Image(d^{n-1}) \arrow[r] & X^n \arrow[r] \arrow[twoheadrightarrow, d] & X^{n+1} \arrow[r] & X^{n+2} \arrow[r] & \cdots \\
		\tau^{\geq n} X \arrow[phantom, r, "" {name=B, yshift=-2ex}] & \cdots \arrow[r] & 0 \arrow[r] & 0 \arrow[r] & \Coker(d^{n-1}) \arrow[r] & X^{n+1} \arrow[r] & X^{n+2} \arrow[r] & \cdots
		\arrow[dash, thick, from=A, to=B]
	\end{tikzcd}\]
	Suku-suku yang dihilangkan sudah jelas, dan dalam arah vertikal hanya
	morfisme selain $\identity$ dan $0$ yang ditandai. Dengan demikian,
	diperoleh morfisme-morfisme di antara kompleks yang tercantum di sebelah
	kiri; dalam penulisan ini, setiap $\Coim(d^i)$ diidentifikasi dengan
	$\Image(d^i)$ melalui isomorfisme pembanding kanonik.
\end{definisi}

% segment-id: o014.li2.chapter3.truncation-functors.p006
Jelas bahwa $\tau^{\leq n}$, $\tilde{\tau}^{\leq n}$,
$\tau^{\geq n}$, dan $\tilde{\tau}^{\geq n}$ semuanya merupakan funktor
aditif. Berlaku
$\tau^{\leq n} = [-n] \circ \tau^{\leq 0} \circ [n]$; hal yang sama berlaku
bagi $\tilde{\tau}^{\leq n}$, $\tau^{\geq n}$, dan
$\tilde{\tau}^{\geq n}$.

% segment-id: o014.li2.chapter3.truncation-functors.p007
Selain itu, jika $m \leq n$, terdapat epimorfisme yang jelas
$\tau^{\geq m} X \twoheadrightarrow \tau^{\geq n} X$ dan
$\tilde{\tau}^{\geq m} X \twoheadrightarrow \tilde{\tau}^{\geq n} X$,
serta monomorfisme yang jelas
$\tau^{\leq m} X \hookrightarrow \tau^{\leq n} X$ dan
$\tilde{\tau}^{\leq m} X \hookrightarrow \tilde{\tau}^{\leq n} X$.

% segment-id: o014.li2.chapter3.truncation-functors.le001
\begin{lema}\label{prop:truncation-ses}
	Untuk setiap $k \in \Z$, morfisme pembanding
	$\tau^{\leq n}X \to \tilde{\tau}^{\leq n}X$ dan
	$\tilde{\tau}^{\geq n}X \to \tau^{\geq n}X$ pada diagram di atas
	menginduksi isomorfisme berikut dalam $\mathcal{A}$\footnote{Catatan
	penerjemah (O014-C051): sumber menyebut ``morfisme-morfisme di atas'', yang
	dapat terbaca sebagai morfisme transisi untuk $m \leq n$ pada paragraf
	tepat sebelumnya; hanya kedua morfisme pembanding yang ditampilkan di sini
	yang menginduksi isomorfisme kohomologi tersebut.}
% segment-id: o014.li2.chapter3.truncation-functors.q003
	\begin{align*}
		\Hm^k\left( \tau^{\leq n} X\right) & \rightiso \Hm^k\left( \tilde{\tau}^{\leq n} X\right) \simeq \begin{cases} \Hm^k(X), & k \leq n \\ 0, & k > n \end{cases}, \\
		\Hm^k\left( \tilde{\tau}^{\geq n} X\right) & \rightiso \Hm^k\left( \tau^{\geq n} X\right) \simeq \begin{cases} \Hm^k(X), & k \geq n \\ 0, & k < n \end{cases},
	\end{align*}
	dan barisan-barisan eksak pendek berikut dalam
	$\cate{C}(\mathcal{A})$:
% segment-id: o014.li2.chapter3.truncation-functors.q004
	\begin{equation*}\begin{gathered}
		0 \to \tilde{\tau}^{\leq n-1} X \to \tau^{\leq n} X \to \Hm^n(X)[-n] \to 0, \\
		0 \to \Hm^n(X)[-n] \to \tau^{\geq n} X \to \tilde{\tau}^{\geq n+1} X \to 0, \\
		0 \to \tau^{\leq n} X \to X \to \tilde{\tau}^{\geq n+1} X \to 0, \\
		0 \to \tilde{\tau}^{\leq n-1} X \to X \to \tau^{\geq n} X \to 0, \\
		0 \to \tau^{\leq n} X \to \tilde{\tau}^{\leq n} X \to \Cone\left(\identity_{\Image(d_X^n) [-n-1]} \right) \to 0,
	\end{gathered}\end{equation*}
	Di sini, $\Hm^n(X)[-n]$ dipahami menurut
	Konvensi~\ref{con:concentrated-cplx}. Semua morfisme tersebut funktorial
	terhadap $X$.
\end{lema}
% segment-id: o014.li2.chapter3.truncation-functors.pf001
\begin{bukti}
	Morfisme $\tilde{\tau}^{\leq n-1} X \to \tau^{\leq n} X$ pada komponen
	berderajat $n$ ialah $\Coim(d^{n-1}) \xrightarrow{\sim}
	\Image(d^{n-1}) \hookrightarrow \Ker(d^n)$,\footnote{Catatan penerjemah
	(O014-C052): sumber langsung menulis $\Image(d^{n-1})$ sebagai suku
	berderajat $n$ dari $\tilde{\tau}^{\leq n-1}X$; suku itu sebenarnya
	$\Coim(d^{n-1})$, yang diidentifikasi secara kanonik dengan citranya.}
	sedangkan pada komponen lain ialah $\identity$. Morfisme
	$\tau^{\leq n} X \to \Hm^n(X)[-n]$ pada komponen berderajat $n$ ialah
% segment-id: o014.li2.chapter3.truncation-functors.q005
	\[ \Ker(d^n) \twoheadrightarrow \Ker(d^n)/\Image(d^{n-1}) = \Hm^n(X). \]
	
	Morfisme $\tau^{\geq n} X \to \tilde{\tau}^{\geq n+1} X$ pada komponen
	berderajat $n$ ialah epimorfisme yang diinduksi oleh $d^n$,
	$\Coker(d^{n-1}) \twoheadrightarrow \Image(d^n)$, sedangkan pada komponen
	lain ialah $\identity$. Morfisme
	$\Hm^n(X)[-n] \to \tau^{\geq n} X$ pada komponen berderajat $n$ ialah
% segment-id: o014.li2.chapter3.truncation-functors.q006
	\[ \Ker(d^n)/\Image(d^{n-1}) \hookrightarrow X^n/\Image(d^{n-1}) = \Coker(d^{n-1}). \]
	
	Verifikasi yang tersisa semuanya rutin dan diserahkan kepada pembaca
	sebagai latihan.
\end{bukti}

% segment-id: o014.li2.chapter3.truncation-functors.p008
Jadi, $\tau^{\leq n}$ dan $\tilde{\tau}^{\leq n}$ (atau
$\tau^{\geq n}$ dan $\tilde{\tau}^{\geq n}$) memang memenggal kohomologi
pada derajat $> n$ (atau $< n$).

% segment-id: o014.li2.chapter3.truncation-functors.re001
\begin{catatan}\label{rem:truncation-sigma}
	Untuk kompleks $X \in \Obj(\cate{C}(\mathcal{A}))$, penerapan funktor
	$\sigma$ dari Definisi--Proposisi~\ref{def:sigma} memberikan
	\[
	\sigma X \in \Obj(\cate{C}(\mathcal{A}^{\opp})) =
	\Obj(\cate{C}(\mathcal{A}^{\opp})^{\opp}).
	\]
	Dengan mengabaikan beberapa
	tanda $\pm$ yang tidak berpengaruh dalam konteks ini, operasi tersebut sama
	dengan membalik arah panah dalam kompleks lalu membalik indeks derajatnya.
	Berdasarkan dualitas antara $\Ker$ dan $\Coker$, di bawah operasi ini
	$\tau^{\geq n}(X)$ bersesuaian dengan
	$\tau^{\leq -n}(\sigma X) \in
	\Obj(\cate{C}(\mathcal{A}^{\opp}))$, dan seterusnya. Demikian pula, karena
	$\Image$ dan $\Coim$ saling dual,
	$\tilde{\tau}^{\geq n}(X) \in \Obj(\cate{C}(\mathcal{A}))$ bersesuaian
	dengan $\tilde{\tau}^{\leq -n}(\sigma X) \in
	\Obj(\cate{C}(\mathcal{A}^{\opp}))$.
\end{catatan}

% segment-id: o014.li2.chapter3.truncation-functors.p009
Dibandingkan dengan $\tilde{\tau}^{\leq n}$ dan
$\tilde{\tau}^{\geq n}$, funktor $\tau^{\leq n}$ dan $\tau^{\geq n}$
mempunyai beberapa sifat yang lebih mudah digunakan. Yang pertama ialah
adjungsi.

% segment-id: o014.li2.chapter3.truncation-functors.pr001
\begin{proposisi}\label{prop:truncation-adjunction}
	Untuk setiap $n \in \Z$, terdapat adjungsi berikut:
% segment-id: o014.li2.chapter3.truncation-functors.q007
	\begin{gather*}
		\Hom_{\cate{C}(\mathcal{A})}(X, Y) \simeq \Hom_{\cate{C}^{\leq n}(\mathcal{A})}\left(X, \tau^{\leq n} Y \right), \quad X \in \Obj(\cate{C}^{\leq n}(\mathcal{A})) \\ \Hom_{\cate{C}(\mathcal{A})}(X, Y) \simeq \Hom_{\cate{C}^{\geq n}(\mathcal{A})}\left(\tau^{\geq n} X, Y \right), \quad Y \in \Obj(\cate{C}^{\geq n}(\mathcal{A})).
	\end{gather*}
\end{proposisi}
% segment-id: o014.li2.chapter3.truncation-functors.pf002
\begin{bukti}
	Hal ini jelas dari definisi pemenggalan. Pembaca dipersilakan menuliskan
	morfisme unit dan kounit yang bersesuaian: morfisme-morfisme tersebut ialah
	morfisme yang ditampilkan dalam Definisi~\ref{def:truncation-cplx} atau
	$\identity$.
\end{bukti}

% segment-id: o014.li2.chapter3.truncation-functors.p010
Selanjutnya, $\tau^{\leq n}$ dan $\tau^{\geq n}$ turun ke tataran
$\cate{K}(\mathcal{A})$.

% segment-id: o014.li2.chapter3.truncation-functors.dp001
\begin{definisiproposisi}\label{def:K-truncation-functors}
	\index[sym1]{taun}
	Untuk setiap $n \in \Z$, funktor $\tau^{\leq n}$ (atau
	$\tau^{\geq n}$) secara alami menginduksi funktor dari
	$\cate{K}(\mathcal{A})$ ke dirinya sendiri, yang tetap dinotasikan dengan
	$\tau^{\leq n}$ (atau $\tau^{\geq n}$).
\end{definisiproposisi}
% segment-id: o014.li2.chapter3.truncation-functors.pf003
\begin{bukti}
	Misalkan $h \in \Hom^{-1}(X, Y)$. Untuk kasus $\tau^{\leq n}$, definisikan
	$\overline{h}^n := h^n \circ \bigl(\Ker(d_X^n) \hookrightarrow X^n\bigr)$.
	Definisikan $\overline{h} \in
	\Hom^{-1}\left( \tau^{\leq n} X, \tau^{\leq n} Y \right)$ menurut
	diagram berikut:
% segment-id: o014.li2.chapter3.truncation-functors.g002
	\[\begin{tikzcd}
		\cdots \arrow[r] & X^{n-2} \arrow[r] \arrow[ld, "h^{n-2}" description] & X^{n-1} \arrow[r] \arrow[ld, "{h^{n-1}}" description] & \Ker(d_X^n) \arrow[r] \arrow[ld, "{\overline{h}^n}" description] & 0 \arrow[r] \arrow[ld] & 0 \arrow[r] \arrow[ld] & \cdots \\
		\cdots \arrow[r] & Y^{n-2} \arrow[r] & Y^{n-1} \arrow[r] & \Ker(d_Y^n) \arrow[r] & 0 \arrow[r] & 0 \arrow[r] & \cdots ;
	\end{tikzcd}\]
	Mudah dilihat bahwa
	$\tau^{\leq n}\left( d^{-1} h \right) = d^{-1} \overline{h}$.
	
	Untuk kasus $\tau^{\geq n}$, sebagai gantinya definisikan
	$\underline{h}^{n+1} := \bigl(Y^n \twoheadrightarrow
	\Coker(d_Y^{n-1})\bigr) \circ h^{n+1}$, lalu gunakan ini untuk
	mendefinisikan $\underline{h} \in
	\Hom^{-1}\left( \tau^{\geq n} X, \tau^{\geq n} Y \right)$; langkah
	selanjutnya serupa. Kedua kasus tersebut tentu saling dual.
\end{bukti}

% segment-id: o014.li2.chapter3.truncation-functors.p011
Sifat sederhana tetapi penting berikut langsung diperoleh dari
Definisi~\ref{def:truncation-cplx}.

% segment-id: o014.li2.chapter3.truncation-functors.pr002
\begin{proposisi}\label{prop:truncate-twice}
	Untuk semua bilangan bulat $a < b$ dan $n$, serta
	$X \in \Obj(\cate{C}(\mathcal{A}))$, berlaku\footnote{Catatan penerjemah
	(O014-C053): sumber menulis $X \in \Obj(\mathcal{A})$, tetapi funktor
	pemenggalan dan $\Hm^n$ pada proposisi ini bekerja pada kompleks.}
% segment-id: o014.li2.chapter3.truncation-functors.q008
	\begin{gather*}
		\tau^{\leq a} \tau^{\geq b}(X) = 0 = \tau^{\geq b} \tau^{\leq a}(X), \\
		\tau^{\leq n }\tau^{\geq n}(X) \simeq \Hm^n(X)[-n] \simeq \tau^{\geq n} \tau^{\leq n}(X) \quad \text{(isomorfisme kanonik).}
	\end{gather*}
\end{proposisi}
